\chapter{Validation and Performance}\label{r02:ch:results}

All numbers in this chapter were measured on the Banana~Pi BPI-F3 between 2026-11-01 and
2026-11-03 (board time) and are recorded in
\file{docs/notes2/rvv-vls256-only-progress-opus5.5.md}; the raw logs are on the board
under \file{/home/bananapi/saleh/}.

\section{Correctness}\label{r02:sec:results-tests}

\paragraph{\ps{} tests.}
\begin{itemize}
  \item Laptop (x86): the three RVV test modules give 77 passed and 7 skipped (the
        hardware tests); the syntax check compiles all eight kernel shapes with the
        riscv64 GCC cross compiler. A regression run over \file{tests/nbackend},
        \file{test_quicktests}, \file{test_type_cast}, \file{test_reduction},
        \file{test_functions} and \file{tests/jit} gives 636 passed, 250 skipped,
        2 expected failures and no failures, so the x86 paths are unaffected.
  \item Board (native GCC 14.2): 84 of 84 RVV tests pass, including bit-exact Philox
        against \code{GenericCPU} and C truncation semantics; the 39 tests of
        \file{test_vectorization.py} pass (all of them run as RVV on the board); and 126
        RVV-parametrized cases of the shared kernel tests (casts, reductions with
        OpenMP, functions, iteration slices, domain kernels, the JIT) pass.
\end{itemize}

\paragraph{waLBerla tests.}
\begin{itemize}
  \item \code{TestSweepGeneration}: 6 of 6 pass (1D and 2D fields, v8 fields, field
        swaps). This test pins its sweeps to \code{Target.CPU} on every platform, so it
        validates the sweepgen machinery and the new build flags, not the RVV backend.
  \item An extra 3D stencil sweep on a \cls{GhostLayerField}, generated with
        vectorization and no target override, includes \file{rvv.hpp}, uses
        \code{vle64}/\code{vfadd}/\code{vfmul}/\code{vfmv}/\code{vse64}, and compiles
        with waLBerla's full warning set (\code{-Wall -Wconversion -Wshadow -pedantic})
        without a single warning.
  \item \code{TestLbmScenarios_CXX}: 9 of 9 pass (fully periodic, mirrored half channel
        and free-slip pipe, each with 1, 2 and 4 MPI processes). This is the default
        waLBerla path on the K1 after the change: \code{CurrentCPU} = \code{RISCV_RVV},
        not vectorized, with lbmpy 2.1.1, boundaries, and the v8 halo exchange over MPI.
\end{itemize}

\paragraph{Checksums of the benchmark.}
Every run prints the total kinetic energy after the last step.
\begin{itemize}
  \item \textbf{Double precision:} all six variants of both compilers agree to all
        twelve printed digits ($3.775155682480\cdot10^{2}$ for $64^3$ cells and 200
        steps), on one core and with 4 and 8 processes. The explicit RVV sweeps compute
        exactly what the scalar ones compute.
  \item \textbf{Single precision:} variants with the same expression structure agree
        exactly (\code{scalar} = \code{autovec}). Between structures, results differ by
        at most $4.7\cdot10^{-6}$ relative, about 60 single-precision ulps after
        200 steps over 262\,144 cells. Global CSE reorders the arithmetic, and the
        compilers contract scalar code into FMAs but not RVV intrinsics, so this is
        expected. The explicit sweeps give \emph{bit-identical} results with GCC and
        Clang: the intrinsics fix the order of operations.
\end{itemize}

\section{Benchmark setup}\label{r02:sec:results-setup}

\code{SweepgenLbmCPU} (\cref{r02:sec:app}): D3Q19 SRT, fully periodic, $64^3$ cells per
process. Single core: 200 timed steps after 10 warm-up steps, pinned to core 2, best of
3 runs (spread at most 1.2\,\% where recorded). Multi-core: 100 steps, \code{mpirun --bind-to core},
best of 2 runs (spread at most 3.9\,\%). Four configurations: \{GCC 14.2, Clang 18.1.8\}
$\times$ \{double, single\} precision, each with the six variants of
\cref{r02:sec:app}.

\section{Single core}\label{r02:sec:results-single}

\begin{table}[htbp]
  \centering\small
  \caption{Single-core throughput in MLUP/s ($64^3$ cells, best of 3). In parentheses:
    speedup over the plain scalar build of the same configuration.}\label{r02:tab:single}
  \begin{tabular}{@{}lllll@{}}
    \toprule
    \textbf{Variant} & \textbf{GCC f64} & \textbf{Clang f64} & \textbf{GCC f32} & \textbf{Clang f32}\\
    \midrule
    scalar        & 0.90 (1.00$\times$) & 0.77 (1.00$\times$) & 1.39 (1.00$\times$) & 1.10 (1.00$\times$)\\
    scalar\_cse   & 0.96 (1.06$\times$) & 0.86 (1.12$\times$) & 1.56 (1.12$\times$) & 1.22 (1.11$\times$)\\
    autovec       & 0.91 (1.00$\times$) & 0.77 (1.00$\times$) & 1.39 (1.00$\times$) & 1.11 (1.01$\times$)\\
    autovec\_cse  & 0.96 (1.06$\times$) & 0.87 (1.12$\times$) & 1.60 (1.15$\times$) & 1.24 (1.12$\times$)\\
    explicit      & 1.14 (1.26$\times$) & \textbf{1.01 (1.31$\times$)} & 2.37 (1.70$\times$) & \textbf{2.11 (1.92$\times$)}\\
    explicit\_cse & \textbf{1.18 (1.31$\times$)} & 0.99 (1.28$\times$) & \textbf{2.40 (1.73$\times$)} & 2.06 (1.87$\times$)\\
    \midrule
    \multicolumn{5}{@{}l}{\textit{Best explicit / best autovec}}\\
                  & \textbf{1.23$\times$} & \textbf{1.16$\times$} & \textbf{1.50$\times$} & \textbf{1.70$\times$}\\
    \bottomrule
  \end{tabular}
\end{table}

\begin{table}[htbp]
  \centering\small
  \caption{Vector instructions in the generated sweep object file (\texttt{llvm-objdump}).}\label{r02:tab:insns}
  \begin{tabular}{@{}lrrrr@{}}
    \toprule
    \textbf{Variant} & \textbf{GCC f64} & \textbf{Clang f64} & \textbf{GCC f32} & \textbf{Clang f32}\\
    \midrule
    scalar        & 0   & 0   & 0   & 0\\
    autovec       & 138 & 55  & 138 & 55\\
    explicit      & 665 & 592 & 665 & 590\\
    explicit\_cse & 608 & 561 & 608 & 561\\
    \bottomrule
  \end{tabular}
\end{table}

\Cref{r02:tab:single,r02:tab:insns} support four conclusions:
\begin{enumerate}
  \item \textbf{Before this work, waLBerla's LBM effectively ran scalar on the K1.}
        Neither compiler's auto-vectorizer does anything useful with the D3Q19 sweep:
        GCC emits 138 vector instructions, Clang 55, and \code{autovec} runs at the
        speed of \code{scalar}. The generic kernel takes the $x$ stride as a run-time
        argument (\cref{r02:sec:generated-sweep}), so the compiler cannot prove that
        consecutive cells are contiguous.
  \item \textbf{The explicit RVV sweep is the fastest build in every configuration}:
        1.16--1.23$\times$ the best auto-vectorized build in double precision,
        1.50--1.70$\times$ in single precision.
  \item \textbf{Single precision gains more.} A 256-bit register holds 8 floats instead
        of 4 doubles, and each cell moves half the bytes.
  \item \textbf{Global CSE helps GCC but not Clang.} With GCC, \code{explicit_cse} is
        3\,\% faster than \code{explicit} in f64; with Clang it is 2--3\,\% slower.
        Clang removes repeated intrinsic calls by itself, GCC does not
        (\cref{r02:sec:backend-glance}). Clang is slower than GCC in every variant
        (scalar f64: 0.77 versus 0.90), so its speedup ratios look better only because
        its baseline is lower. The fastest build overall is GCC with \code{explicit_cse}.
\end{enumerate}

\section{Multiple cores}\label{r02:sec:results-multi}

\begin{table}[htbp]
  \centering\small
  \caption{Throughput with 4 and 8 MPI processes in MLUP/s ($64^3$ cells per process,
    best of 2), and the speedup over scalar.}\label{r02:tab:multi}
  \begin{tabular}{@{}llllll@{}}
    \toprule
    \textbf{Processes} & \textbf{Variant} & \textbf{GCC f64} & \textbf{Clang f64} & \textbf{GCC f32} & \textbf{Clang f32}\\
    \midrule
    4 & scalar        & 3.30 (1.00$\times$) & 2.82 (1.00$\times$) & 5.25 (1.00$\times$) & 4.31 (1.00$\times$)\\
      & autovec\_cse  & 3.48 (1.05$\times$) & 3.12 (1.11$\times$) & 5.57 (1.06$\times$) & 4.81 (1.12$\times$)\\
      & explicit      & 3.84 (1.16$\times$) & 3.56 (1.26$\times$) & 7.65 (1.46$\times$) & 7.13 (1.66$\times$)\\
      & explicit\_cse & 3.89 (1.18$\times$) & 3.51 (1.24$\times$) & 7.81 (1.49$\times$) & 6.96 (1.62$\times$)\\
    \midrule
    8 & scalar        & 4.90 (1.00$\times$) & 4.48 (1.00$\times$) & 8.55 (1.00$\times$) & 7.50 (1.00$\times$)\\
      & autovec\_cse  & 5.00 (1.02$\times$) & 4.78 (1.07$\times$) & 9.00 (1.05$\times$) & 8.07 (1.08$\times$)\\
      & explicit      & 5.24 (1.07$\times$) & 5.07 (1.13$\times$) & 10.38 (1.21$\times$) & 10.13 (1.35$\times$)\\
      & explicit\_cse & 5.26 (1.07$\times$) & 5.03 (1.12$\times$) & 10.45 (1.22$\times$) & 9.95 (1.33$\times$)\\
    \bottomrule
  \end{tabular}
\end{table}

\begin{figure}[htbp]
  \centering
  % Best explicit RVV build over the best auto-vectorized build, per configuration and
% number of MPI processes (data: Table "best explicit / best autovec", Chapter 11).
% Sequential encoding: more processes = darker blue.
\begin{tikzpicture}[x=1cm,y=2.6cm,font=\footnotesize]
  \def\bw{0.62}   % bar width
  \def\gap{0.07}  % surface gap between bars
  % recessive grid and axis
  \foreach \v in {0.5,1.5} \draw[psgray!25,line width=0.3pt] (-0.4,\v) -- (11.9,\v);
  \foreach \v in {0,0.5,1,1.5}
    \node[anchor=east,font=\scriptsize,text=psgray!90!black] at (-0.45,\v) {\pgfmathprintnumber[fixed,precision=1,min exponent for 1000 sep=5]{\v}$\times$};
  \draw[psgray!80,line width=0.4pt] (-0.4,0) -- (11.9,0);
  % data: group / label / np1 / np4 / np8
  \foreach \g/\lab/\a/\b/\c in {0/{GCC, f64}/1.23/1.12/1.05, 1/{Clang, f64}/1.16/1.14/1.06,
                                 2/{GCC, f32}/1.50/1.40/1.16, 3/{Clang, f32}/1.70/1.48/1.26} {
    \pgfmathsetmacro{\xo}{\g*3.1+0.35}
    \fill[psblue!30,rounded corners=0.8pt] (\xo,0) rectangle ++(\bw,\a);
    \fill[psblue!62,rounded corners=0.8pt] (\xo+\bw+\gap,0) rectangle ++(\bw,\b);
    \fill[psblue!100!black,rounded corners=0.8pt] (\xo+2*\bw+2*\gap,0) rectangle ++(\bw,\c);
    \node[anchor=north,font=\footnotesize] at (\xo+1.5*\bw+\gap,-0.04) {\lab};
  }
  % baseline: best auto-vectorized build of the same configuration
  \draw[psred!80!black,dashed,line width=0.7pt] (-0.4,1) -- (11.9,1);
  \node[anchor=west,font=\scriptsize,text=psred!80!black,align=left] at (11.95,1.0) {best\\autovec};
  % legend
  \begin{scope}[shift={(0.35,1.8)}]
    \foreach \x/\c/\t in {0/psblue!30/1 process, 2.4/psblue!62/4 processes (one cluster),
                          7.0/psblue!100!black/8 processes (whole chip)} {
      \fill[\c,rounded corners=0.8pt] (\x,0) rectangle ++(0.32,0.09);
      \node[anchor=west,font=\scriptsize] at (\x+0.38,0.045) {\t};
    }
  \end{scope}
\end{tikzpicture}

  \caption[Speedup of explicit RVV over auto-vectorization]{Best explicit RVV build
    divided by the best auto-vectorized build of the same configuration, with 1, 4 and 8
    MPI processes. The gain shrinks as more cores share the memory system.}
  \label{r02:fig:speedup}
\end{figure}

The explicit sweeps stay ahead with more cores, but the margin shrinks
(\cref{r02:fig:speedup}): with all eight cores it is 1.05--1.06$\times$ in double
precision and 1.16--1.26$\times$ in single precision. Going from 4 to 8 processes,
i.e.\ adding the second cluster, raises the throughput of the explicit GCC f64 build only
from 3.89 to 5.26 MLUP/s (1.35$\times$).

\section{Why the gain shrinks}\label{r02:sec:results-bandwidth}

The generated \code{StreamCollide} reads 19 PDFs and writes 19 PDFs, $\rho$ and three
velocity components: at least 42 values per cell, 336 bytes in double precision and 168
in single precision (write-allocate traffic not counted). Multiplying by the
throughput gives the memory traffic:

\begin{center}\small
\begin{tabular}{@{}lrrrrrr@{}}
  \toprule
  & \multicolumn{2}{c}{\textbf{explicit\_cse, GCC}} & \multicolumn{2}{c}{\textbf{explicit, Clang}} & \multicolumn{2}{c}{\textbf{scalar, GCC}}\\
  \cmidrule(lr){2-3}\cmidrule(lr){4-5}\cmidrule(l){6-7}
  \textbf{Processes} & f64 & f32 & f64 & f32 & f64 & f32\\
  \midrule
  1 & 396 & 403 & 339 & 354 & 302 & 234\\
  4 & 1307 & 1312 & 1196 & 1198 & 1109 & 882\\
  8 & 1767 & 1756 & 1704 & 1702 & 1646 & 1436\\
  \bottomrule
\end{tabular}

\smallskip
{\footnotesize Lower bound of the memory traffic in MB/s (MLUP/s $\times$ bytes per cell).}
\end{center}

The explicit sweeps move the \emph{same} number of bytes per second in double and single
precision: 396 versus 403\,MB/s on one core, 1307 versus 1312 on four, 1767 versus 1756 on
eight. That is the signature of a bandwidth limit. Single precision runs twice as many
cells per second only because each cell is half the bytes. The scalar builds stay below
the limit, so vectorization helps them; the more cores share the memory system, the
closer even scalar code gets to the limit (1646 of about 1770\,MB/s with eight cores in
f64), and the less there is to gain. On one core, about 400\,MB/s is a per-core limit,
not the DRAM limit: the in-order X60 core can keep only a few cache misses in flight.

\begin{keyidea}[What the RVV support achieves]
  The backend removes the compute bottleneck of the waLBerla LBM sweep on the K1. Beyond
  that, the sweep is limited by memory bandwidth, which no code generator can fix.
  Further gains need less traffic per cell: in-place streaming (AA pattern, esoteric
  twist), single-precision storage, or fusing more work into one pass.
\end{keyidea}

\section{Context: kernel-level experiments}\label{r02:sec:results-context}

Before the waLBerla runs, the same backend was studied on a D2Q9 LBM harness
(\file{python/learn-lb-pystencil/lb-pystencils.py}) and on standalone kernels
(\file{misc/rvv-vs-autovec/}). Their findings explain the waLBerla numbers:
\begin{itemize}
  \item With the inner stride a run-time value, as in waLBerla and the \ps{} JIT,
        GCC's auto-vectorizer turns every field access into indexed gathers and scatters
        (\code{vluxei64}/\code{vsuxei64}), while the explicit backend uses unit-stride
        loads. With GCC, the explicit kernel was 1.24--1.31$\times$ (f64) and
        1.35--1.53$\times$ (f32) faster than the best auto-vectorized build.
  \item Clang versions the loop on stride $=1$ and auto-vectorizes at LMUL 2, so for
        Clang the stride contract gives no advantage. The weakest combination is Clang
        with explicit intrinsics: it keeps all constant splats in vector registers and
        spills, and schedules poorly for the in-order X60. Rewriting splats as
        vector--scalar \code{.vf} operations (\code{vfmul_vf} instead of
        \code{vfmul_vv} with a \code{vfmv} splat) removed the spills and gained 7--25\,\%
        with Clang and 1.5--13\,\% with GCC.
  \item Hand-inserted FMAs gave no gain with GCC, so the backend does not emit them.
\end{itemize}
